\documentclass[10pt,a4paper]{amsart}
\usepackage[margin=19mm]{geometry}
\usepackage{amsmath,amssymb,mathtools}
\usepackage[scr=rsfs,cal=euler]{mathalfa}
\usepackage{xfrac}
\usepackage[unicode,hidelinks,bookmarks=false]{hyperref}
\newcommand{\ie}{i.e.\ }
\newcommand{\cf}{cf.\ }
\newcommand{\resp}{respectively\ }
\newcommand{\Subsection}{\subsection}
\providecommand{\nobreakdash}{\nobreak\hbox{-}}
\setlength{\emergencystretch}{2em}
\multlinegap0pt
\usepackage{amssymb}
\usepackage{mathtools}
\usepackage[shortlabels]{enumitem}
\newtheorem{theorem}{Theorem}
\newcommand{\Z}{\mathbb{Z}}
\title{The Golden Sieve}
\author{Benoit Cloitre}
\date{Source: arXiv:2602.17735v2 (CC BY 4.0)}
\begin{document}
\maketitle

\noindent\emph{Source and licence.} The Golden Sieve. Version arXiv:2602.17735v2 (CC BY 4.0), distributed by arXiv under the Creative Commons Attribution 4.0 International licence, http://creativecommons.org/licenses/by/4.0/, as recorded in the arXiv licence metadata for this exact version and shown on the abstract page https://arxiv.org/abs/2602.17735v2. Full licence notice: \texttt{licenses/arxiv-2602.17735-CC-BY-4.0.txt}.\par\smallskip

\noindent\textit{Editorial scope.} Counted unit: the article's theorem thm:metallic-beatty (source0.tex lines 1794--1816). The article's complete original proof of it is delivered with it (source0.tex lines 1818--1889, one \texttt{proof} environment of 72 lines). No mathematics is written, repaired or re-derived: the statement and the proof are the article's own lines, in the article's own notation, and no retained line is substituted or re-flowed.  No further source-local context is needed: every cross-reference of every retained line resolves inside the two delivered spans.  Nothing was omitted from any retained span, so the package carries no omission record.  No retained line contains a citation, so the wrapper carries no bibliography and no bibliography entry.  No retained line contains a figure environment, an image-inclusion command or drawing code, so no image is delivered and the package declares no asset dependency.  Editorial formatting only, in addition: an AMS \texttt{amsart} preamble that restates the article's own declarations and macros used by the retained lines, namely the environments \texttt{theorem} on separate counters, together with the standard packages those lines need.  The retained environments print in this document as Theorem 1 in order of appearance, whereas the article prints the same items with the numbering of its own full document.  The complete native source of the article as distributed in the arXiv e-print for this exact version is bundled unchanged as \texttt{sources/194983/source0.tex}.
\begin{theorem}\label{thm:metallic-beatty}
Let $k\ge1$.
\begin{enumerate}
\item[\textup{(a)}] \textbf{Case $y=k+1$, $z=k$}
  \textup{(large gap on hit).}
  Let $\alpha=M_k:=(k+\sqrt{k^2+4})/2$, the positive root of
  $\alpha^2=k\alpha+1$, and define
  \begin{equation}\label{eq:metallic-beta}
    \beta_k=-\frac{(k-1)\,\alpha}{\alpha+1}.
  \end{equation}
  Then $s_n=\lfloor M_k\,n+\beta_k\rfloor$ for all $n\ge1$.

\item[\textup{(b)}] \textbf{Case $y=k$, $z=k+1$}
  \textup{(large gap on miss).}
  Let $\alpha=R_k:=\bigl((k+1)+\sqrt{(k+1)^2-4}\bigr)/2$, the
  positive root of $\alpha^2=(k+1)\alpha-1$, and define
  \begin{equation}\label{eq:reverse-beta}
    \beta_k^{\mathrm{r}}=-\frac{(k-1)\,\alpha}{\alpha-1}.
  \end{equation}
  Then $s_n=\lfloor R_k\,n+\beta_k^{\mathrm{r}}\rfloor$ for
  all $n\ge1$.
\end{enumerate}
\end{theorem}

\begin{proof}[Proof of~\textup{(a)}]
Set $b(n)=\lfloor \alpha n+\beta\rfloor$, where $\alpha=M_k$ and
$\beta=\beta_k$, and let $B=\operatorname{Im}(b)$.

We first check that $b(1)=1$.  Since
\[
\alpha+\beta
=\alpha\Bigl(1-\frac{k-1}{\alpha+1}\Bigr)
=\frac{\alpha(\alpha+2-k)}{\alpha+1},
\]
and $\alpha^2=k\alpha+1$, this becomes
\[
\frac{1+2\alpha}{\alpha+1}
=2-\frac{1}{\alpha+1},
\]
which lies in $(1,2)$.  Hence $b(1)=1$.

Now $k<\alpha<k+1$, so the gap $b(n+1)-b(n)$ is either $k$ or $k+1$.
More precisely,
\[
b(n+1)-b(n)=k+1
\iff \{\alpha n+\beta\}\ge (k+1)-\alpha
=\frac{\alpha-1}{\alpha}.
\]
On the other hand, a standard inversion argument shows that
\[
m\in B
\iff \Bigl\{\frac{m-\beta}{\alpha}\Bigr\}\ge 1-\frac1\alpha
=\frac{\alpha-1}{\alpha}.
\]
So the same threshold appears on both sides.

It remains to relate the two fractional parts.  We compute
\[
\alpha n+\beta-\frac{n-\beta}{\alpha}
=n\Bigl(\alpha-\frac1\alpha\Bigr)
+\beta\Bigl(1+\frac1\alpha\Bigr).
\]
Since $1/\alpha=\alpha-k$, we have $\alpha-1/\alpha=k$, and by the
definition of $\beta$,
\[
\beta\Bigl(1+\frac1\alpha\Bigr)
=\beta\,\frac{\alpha+1}{\alpha}
=-(k-1).
\]
Therefore
\[
\alpha n+\beta-\frac{n-\beta}{\alpha}=kn-(k-1)\in\Z.
\]
Because $\alpha$ is irrational and $\alpha+\beta\notin\Z$, the quantity
$\alpha n+\beta$ is never an integer.  Hence
\[
\{\alpha n+\beta\}
=
\Bigl\{\frac{n-\beta}{\alpha}\Bigr\}
\qquad (n\ge1).
\]

We may now combine the previous observations.  The gap $b(n+1)-b(n)$
equals $k+1$ exactly when
\[
\{\alpha n+\beta\}\ge \frac{\alpha-1}{\alpha},
\]
which is equivalent to
\[
\Bigl\{\frac{n-\beta}{\alpha}\Bigr\}\ge \frac{\alpha-1}{\alpha},
\]
and therefore to $n\in B$.  This is exactly the defining rule of the
$(1,1,k{+}1,k)$-hiccup sequence.  The two sequences have the same
initial value and the same rule for choosing the next gap, so they
coincide for all $n\ge1$.
\end{proof}

\end{document}
